% ---------------------------------------------------------------------------
%  08-evaluation.tex -- what was measured, and what it does and does not show
%
%  PROVENANCE.  Every number in this section comes from
%  benchmarks/clean_results/ -- the cleaned grid, in which the 600 s runs have
%  had the 2 h re-runs substituted in where those exist (see CLEANING_REPORT).
%  The raw benchmarks/results/ tree is NOT used.  All PTCTLMC rows are the
%  chunk-8 configuration (--omp-chunk 8), best of two repetitions.
% ---------------------------------------------------------------------------
\section{What was measured}
\label{sec:eval}

All numbers below are from \file{benchmarks/clean\_results/}, the cleaned grid
in which the $600$\,s runs have had the $2$\,h re-runs substituted in where
those exist.  The grid is five families (\code{phil}, \code{fischer},
\code{twolock}, \code{viking}, \code{single}) $\times$ $\{16,24,32\}$ worker
threads $\times$ chunk sizes $\{2,4,8,16\}$, two repetitions each, against
\code{verifyta} (\uppaal{}~$5.0.0$) run single-threaded as the oracle.  Unless
stated otherwise the PTCTLMC column is the \emph{chunk-8} configuration, best
of the two repetitions.

\subsection{Speedup against \code{verifyta}, at all three thread counts}

The ratio columns mix two effects -- parallelism and a different search -- and
must be read together with the state counts in \cref{sec:eval-states}.

\begin{center}
\small
\begin{tabular}{@{}lrrrrrrr@{}}
\toprule
& \code{verifyta} & \multicolumn{3}{c}{\tool{} wall-clock (s)}
& \multicolumn{3}{c}{speedup} \\
\cmidrule(lr){3-5}\cmidrule(lr){6-8}
model & (s) & $16$ thr & $24$ thr & $32$ thr & $\times16$ & $\times24$ & $\times32$ \\
\midrule
\code{phil11}    & $4.42$    & $0.61$   & $0.45$   & $0.40$   & $7.2$ & $9.9$  & $\mathbf{11.1}$ \\
\code{phil12}    & $19.66$   & $2.66$   & $1.97$   & $1.58$   & $7.4$ & $10.0$ & $\mathbf{12.4}$ \\
\code{phil13}    & $96.60$   & $16.97$  & $11.84$  & $9.42$   & $5.7$ & $8.2$  & $\mathbf{10.3}$ \\
\code{phil14}    & $533.16$  & $117.91$ & $86.00$  & $68.13$  & $4.5$ & $6.2$  & $\mathbf{7.8}$  \\
\addlinespace[2pt]
\code{viking15}  & $18.26$   & $2.67$   & $1.98$   & $1.70$   & $6.9$ & $9.2$  & $\mathbf{10.8}$ \\
\code{viking17}  & $97.86$   & $13.21$  & $9.77$   & $8.26$   & $7.4$ & $10.0$ & $\mathbf{11.8}$ \\
\code{viking20}  & $1253.98$ & $159.21$ & $114.36$ & $90.98$  & $7.9$ & $11.0$ & $\mathbf{13.8}$ \\
\addlinespace[2pt]
\code{fischer12} & $91.39$   & $23.62$  & $17.49$  & $14.52$  & $3.9$ & $5.2$  & $\mathbf{6.3}$  \\
\code{fischer14} & $1289.10$ & $371.74$ & $275.75$ & $230.10$ & $3.5$ & $4.7$  & $\mathbf{5.6}$  \\
\addlinespace[2pt]
\code{twolock13} & $358.18$  & $90.96$  & $66.65$  & $54.36$  & $3.9$ & $5.4$  & $\mathbf{6.6}$  \\
\code{twolock14} & $1310.24$ & $361.59$ & $268.67$ & $217.79$ & $3.6$ & $4.9$  & $\mathbf{6.0}$  \\
\addlinespace[2pt]
\code{single6}   & $4.79$    & $3.07$   & $3.61$   & $3.31$   & $1.6$ & $1.3$  & $1.4$ \\
\code{viking4}   & $0.06$    & $0.34$   & $0.07$   & $0.08$   & $0.2$ & $0.9$  & $0.8$ \\
\bottomrule
\end{tabular}
\end{center}

At $32$ threads the engine is $5.6$--$13.8\times$ \code{verifyta} on the four
large families.  The two rows at the bottom are the ones that do not fit the
story and are reported for that reason: \code{single6} barely gains because its
entire passed set is one discrete state (see below), and \code{viking4}
finishes in $60$\,ms, where thread-team creation and the interner's initial
allocations dominate everything else.

\subsection{Scaling within the grid}

\FigScaling

\begin{center}
\small
\begin{tabular}{@{}lrrrrr@{}}
\toprule
model & $16$ thr & $24$ thr & $32$ thr & $16{\to}24$ & $16{\to}32$ \\
\midrule
\code{viking20}  & $159.21$\,s & $114.36$\,s & $90.98$\,s  & $1.39\times$ & $1.75\times$ \\
\code{phil14}    & $117.91$\,s & $86.00$\,s  & $68.13$\,s  & $1.37\times$ & $1.73\times$ \\
\code{twolock14} & $361.59$\,s & $268.67$\,s & $217.79$\,s & $1.35\times$ & $1.66\times$ \\
\code{fischer14} & $371.74$\,s & $275.75$\,s & $230.10$\,s & $1.35\times$ & $1.62\times$ \\
\code{single6}   & $3.07$\,s   & $3.61$\,s   & $3.31$\,s   & $0.85\times$ & $0.93\times$ \\
\bottomrule
\end{tabular}
\end{center}

Doubling the workers from $16$ to $32$ buys $1.62$--$1.75\times$, i.e.\
$81$--$87\%$ parallel efficiency \emph{in that range}.  \code{single6} is the
designed counterexample and is reported deliberately: one automaton, one
location, $N$ independent clocks, therefore \emph{one} discrete state,
therefore one stripe and one antichain.  Every worker contends for the same
lock and the search runs single-file.  The hot-chain dense mirror and the
shared-lock probe were both built for this shape and both help the constant
factor -- before the probe covered mirrored chains, \code{single6} ran at
$733\%$ CPU with \emph{zero} speedup -- but they cannot manufacture parallelism
that the model does not contain.  Any honest evaluation must include a family
like this.

\subsection{Stored state counts}
\label{sec:eval-states}

State counts separate the algorithmic effect from the parallel one, so they
carry more weight than the ratio table.  Where they vary run-to-run
(\cref{sec:eval-nondet}) the observed range is given.

\begin{center}
\small
\begin{tabular}{@{}lrrl@{}}
\toprule
model & \tool{} stored & \code{verifyta} stored & reading \\
\midrule
\code{phil12}    & $802\,871$--$805\,452$      & $1\,140\,361$   & $\approx 30\%$ fewer \\
\code{phil14}    & $10\,613\,715$--$10\,652\,721$ & $14\,833\,412$ & $\approx 28\%$ fewer \\
\code{fischer12} & $4\,758\,404$               & $4\,758\,404$   & identical \\
\code{fischer14} & $49\,309\,338$              & $49\,309\,338$  & identical \\
\code{twolock13} & $15\,465\,049$              & $15\,465\,049$  & identical \\
\code{twolock14} & $49\,538\,714$              & $49\,538\,714$  & identical \\
\code{viking17}  & $13\,631\,486$              & $13\,631\,486$  & identical \\
\code{viking20}  & $143\,654\,910$             & $143\,654\,910$ & identical \\
\code{single6}   & $5\,994$--$6\,660$          & $3\,673$        & \tool{} stores \emph{more} \\
\bottomrule
\end{tabular}
\end{center}

Three readings, all of which belong in a paper:

\begin{enumerate}[leftmargin=1.5em, itemsep=3pt]
\item \textbf{Where the state counts are identical} (\code{viking},
      \code{fischer}, \code{twolock}), the two tools are exploring the same
      symbolic object, so the $5.6$--$13.8\times$ at $32$ threads is
      attributable to parallelism and to the storage engineering of
      \cref{sec:storage} -- not to a different search.  This is the cleanest
      comparison available and is the one to lead with.

\item \textbf{Where they differ in \tool{}'s favour} (\code{phil}), part of the
      gain is algorithmic: $\approx 30\%$ fewer stored states from the
      convex-union merge.

\item \textbf{Where \tool{} loses} (\code{single6}, $63$--$81\%$ \emph{more}
      states), the model is the one whose antichain is a single long chain of
      mutually incomparable zones, where widening is the right mechanism and
      pruning is not.  This is exactly the case the adaptive switch
      (\cref{thm:adapt}) exists for; with it on, the stored count falls to
      $3676$ against \uppaal{}'s $3673$.
\end{enumerate}

\subsection{State-count nondeterminism, measured}
\label{sec:eval-nondet}

\cref{thm:determinism} predicts that the verdict is schedule-independent but
the stored antichain is not.  The grid confirms it, and does so in a way that
is more informative than the theorem alone:

\begin{center}
\small
\begin{tabular}{@{}llr@{}}
\toprule
family & stored count across the $6$ (threads $\times$ rep) chunk-8 runs & spread \\
\midrule
\code{phil8}--\code{phil14} & $6$ distinct values each & $0.26$--$1.74\%$ \\
\code{single4}--\code{single6} & $4$--$6$ distinct values each & $11.1$--$19.9\%$ \\
\code{fischer}, \code{twolock}, \code{viking} (all sizes) & one value & $0\%$ \\
\bottomrule
\end{tabular}
\end{center}

$10$ of the $43$ completing models vary; $33$ are bit-stable.  The split is not
arbitrary: the families that vary are exactly the families where the
convex-union merge fires, and the families that are stable are exactly those
whose counts match \code{verifyta} to the digit.  This is the predicted
signature -- step~(4) of \code{add} is order-sensitive, so where it never fires
the result is deterministic, and where it fires often the antichain depends on
which worker reached the stripe first.  The verdicts agree in every run.

\subsection{Chunk size is not a neutral knob on every model}

The scaling and comparison tables above fix the chunk at $8$.  That choice is
not neutral, and the grid says so.  At $32$ threads:

\begin{center}
\small
\begin{tabular}{@{}lrrrr@{}}
\toprule
model & \code{-c 2} & \code{-c 4} & \code{-c 8} & \code{-c 16} \\
\midrule
\code{phil14}    & $117.55$\,s & $91.34$\,s  & $68.13$\,s  & $\mathbf{52.62}$\,s \\
\code{viking20}  & $93.61$\,s  & $92.04$\,s  & $90.98$\,s  & $92.77$\,s \\
\code{twolock14} & $217.40$\,s & $217.09$\,s & $217.79$\,s & $217.08$\,s \\
\code{fischer14} & $229.32$\,s & $228.71$\,s & $230.10$\,s & $230.95$\,s \\
\code{single6}   & $3.36$\,s   & $3.37$\,s   & $3.31$\,s   & $3.04$\,s \\
\bottomrule
\end{tabular}
\end{center}

Three of the five families are flat to within noise, so chunk $8$ costs them
nothing.  \code{phil14} is not: it runs $2.23\times$ faster at chunk $16$ than
at chunk $2$, and $1.29\times$ faster at chunk $16$ than at the chunk $8$ used
in every table above.  \emph{The reported \code{phil} numbers are therefore
conservative} -- at chunk $16$, \code{phil14} is $10.1\times$ \code{verifyta}
rather than $7.8\times$.  The mechanism is per-item overhead: \code{phil}'s
successors are cheap and numerous, so a small chunk pays the dynamic-schedule
bookkeeping too often, while models with expensive successors are insensitive.
A paper should either report the best chunk per model (stating that it does so)
or report chunk $8$ throughout and note this, but must not mix the two.

\subsection{Failures in the grid, disclosed}

Not every cell completed, and an artifact reviewer will read the CSVs:

\begin{center}
\small
\begin{tabular}{@{}lll@{}}
\toprule
model & outcome & note \\
\midrule
\code{phil15}, \code{single7} & timeout at $600$\,s & complete in the $2$\,h re-runs at chunk $2$ \\
\code{fischer15}              & timeout at $600$\,s & no $2$\,h re-run yet \\
\code{twolock15}              & timeout at $16$ thr & completes for \code{verifyta} in $4926$\,s \\
\code{twolock15/16}           & \textbf{ERR} at $24$ and $32$ thr & aborts at $16.5$ / $21.0$\,GB \\
\bottomrule
\end{tabular}
\end{center}

The \code{ERR} rows deserve a sentence rather than a footnote.  They abort at a
\emph{repeatable} resident size ($16.51$--$16.54$\,GB for \code{twolock15},
$20.93$--$21.10$\,GB for \code{twolock16}), identical at $24$ and $32$ threads,
which is the signature of a capacity limit reached at a fixed zone
\emph{count} rather than of a race or of the machine running out of memory.
The trigger is the zone interner's chunk-table ceiling; because the overflow is
raised as an exception from inside an \code{omp parallel} region -- where an
exception may not leave the structured block -- every worker calls
\code{std::terminate} and the process takes \code{SIGABRT}, so the harness
records \code{ERR} at a nonzero exit rather than a clean out-of-memory message.
The constants have since been resized so that RAM is the binding constraint and
the harness reports an honest \code{MEMOUT}; \textbf{these grid cells predate
that change and should be re-run before submission}, both to confirm the fix
and because leaving unexplained \code{ERR} cells in a published artifact is
worse than reporting a limit.

\subsection{Where the time goes}

The engine has a built-in phase profiler (\code{PTCTLMC\_PHPROF}) reporting the
split between \textsc{expand} and the three post-barrier phases, plus
per-thread busy time (hence tail idle).  Each design decision in
\cref{sec:parallel} was taken in response to a specific number from it:

\begin{center}
\begin{tabular}{@{}p{0.40\linewidth}p{0.53\linewidth}@{}}
\toprule
\textbf{Symptom measured} & \textbf{Change it forced} \\
\midrule
$374\%$ CPU of a possible $800\%$, $615$k voluntary context switches
  & shared-lock probe (\cref{fig:protocol}) \\
$378$ of $759$ CPU-s blocked on \code{std::shared\_mutex}
  & spin-then-yield \code{SpinRW} \\
$36.3$\,s unpacking to feed $1.6$\,s of comparisons
  & minimal-graph inclusion filter (\cref{lem:mingraph}) \\
Threads waiting $20\times$ longer than the interner lock was held
  & $64$-way interner sharding \\
$31.2$ probes per lookup after sharding (was $1.68$)
  & high-bit shard selector (\cref{prop:shard}) \\
$2.8$\,s of a $9.3$\,s run in GC ($30\%$)
  & growth-triggered sweep (\cref{prop:gc}) \\
$0.149$\,s serial union per level, growing with level width
  & bucketed collection (\cref{lem:bucket}) \\
$54\%$ of the arena dead, walked on every probe
  & per-stripe compaction floor \\
$4.6$k convex-union tests per insertion on \code{single6}
  & bounded merge window \\
\bottomrule
\end{tabular}
\end{center}

\subsection{Correctness evidence}

\begin{itemize}[leftmargin=1.2em, itemsep=2pt]
\item \textbf{Differential testing against \code{verifyta}} on every model and
      query in the suite; verdicts agree, and on three of five families the
      stored state counts agree exactly.
\item \textbf{\code{seq} $\equiv$ \code{omp}} is unit-tested; at one thread the
      OpenMP engine takes the \code{single} path and reproduces the sequential
      insertion order bit-for-bit.
\item \textbf{Compiled-in invariant checkers} (built with
      \code{PTCTLMC\_DEBUG\_KNOBS}) that a paper can cite as machine-checked
      evidence for specific lemmas:
      \code{CHECKINTERN} (minimal-graph round-trip is exact),
      \code{CHECKMERGE} (every accepted merge verified by exact federation
      subtraction --- i.e.\ \cref{thm:hull} checked at runtime),
      \code{CHECKINV} (every stored zone satisfies the invariant of every
      process's current location),
      \code{CHECKGC} (a frontier entry's interned record is byte-identical when
      the level expands it --- i.e.\ \cref{prop:gc-cache} checked at runtime).
\end{itemize}

\subsection{What is not measured, and should be before submission}

\begin{itemize}[leftmargin=1.2em, itemsep=2pt]
\item \textbf{No low-thread baseline anywhere in the grid.}  The cleaned
      results contain $16$, $24$ and $32$ threads only -- there is no $1$, $2$,
      $4$ or $8$ thread data in \file{clean\_results/} (nor in the raw tree).
      So absolute speedup versus the sequential engine is \emph{not derivable};
      only the $16 \to 32$ slope is.  A $1/2/4/8$ sweep on the same machine is
      the single most valuable missing measurement, and without it the paper
      cannot state a parallel speedup at all -- only a ratio against
      \code{verifyta}, which conflates parallelism with the algorithm.
\item \textbf{The \code{ERR} cells have not been re-run} since the interner
      ceiling was raised (see above).
\item \textbf{No ablation series.}  The mechanisms have env-var switches
      (\code{PTCTLMC\_NOMERGE}, \code{NOGC}, \code{EXACTSUB}, \code{HOTLEN},
      \code{MERGESCAN}, \code{ADAPT}), so a one-mechanism-at-a-time table is
      cheap to produce and would carry far more weight than the anecdotal
      before/after numbers in the table above.
\item \textbf{RSS is recorded but unused.}  \file{clean\_results/} carries an
      \code{rss\_gb} column for both tools; a memory comparison against
      \code{verifyta} is available at no extra experimental cost.
\end{itemize}
